\section{Open questions}
\label{sec:discussion}

The main results are the two-sided bound $(\sqrt2-o(1))\sqrt n\le k^*(n)\le(2\sqrt{\pi e}+o(1))\sqrt{n\log n}$ for the permutation threshold, with exact values for $n\le40$ (computer-assisted for the upper bounds when $18\le n\le40$), and exact second moments of the line-sum energies of a random arrangement.
The following questions remain.

\begin{enumerate}\setlength\itemsep{2pt}
\item \emph{The order of $k^*(n)$.} Is $k^*(n)\asymp\sqrt n$?
The upper bound uses only that $\gamma$ is a nonzero integer vector with $\|\gamma\|_1\le n^4/4$.
In that class the pigeonhole principle produces, for all large $n$, vectors vanishing to order $\ge\lfloor2\sqrt n\rfloor$ at $1$ (\Cref{sec:kstar_upper}), so a bound of smaller order than $\sqrt n$ needs more information about $\gamma$.
Whether the order $\sqrt{n\log n}$ of \Cref{lem:jensen} is attained within this class is, as far as we know, open; this is separate from realisability, i.e.\ which such vectors are column deviations of arrangements.
For $14\le n\le40$, $k^*(n)$ is the least $k$ for which $\Lambda_{n,k}$ contains no nonzero vector satisfying \eqref{eq:col_feasible} (\Cref{sec:kstar_exact}).  Is this true for all $n$, and is $k^*(n)$ nondecreasing?
\item \emph{The size of $L_n$.} Is $|L_n|\sim N!\,(2\pi)^{-2}\det(\Sigma_U)^{-1/2}\,c_n$ with $c_n$ bounded away from $0$ and $\infty$, i.e.\ does a local limit theorem hold for the four statistics $(U_1,\dots,U_4)$ of \Cref{sec:lowmode_n4}?  The ratio $c_n$ is $0.59$, $0.80$ for $n=3,4$ and about $0.76$, $0.78$ for $n=5,6$.
\item \emph{Other families.} The orthogonality formulation applies to pandiagonal and Franklin squares \cite{loly2007franklin}, with $W$ replaced by the span of the relevant line indicators.  The corresponding thresholds and random energies have not been worked out.
\end{enumerate}
